\section{Problem Setup}
\label{sec:setup}
%======================================================================

\subsection{Environment and agents}

We study a partially observable, adversarial multi-robot setting instantiated
as a nominal $2$-vs-$2$ naval Capture-the-Flag task on a symmetric maritime
arena. Blue is the learned multi-robot team; Red is a scripted or frozen
adversarial opponent. Each side controls $|B|=|R|=2$ autonomous surface vessels.
We model the interaction as a finite-horizon partially observable Markov game
\[
G =
\left(
\mathcal{N},\mathcal{S},\{\mathcal{A}_i\}_{i\in\mathcal{N}},P,
\{r_i\}_{i\in\mathcal{N}},\{\Omega_i\}_{i\in\mathcal{N}},\{O_i\}_{i\in\mathcal{N}},\gamma
\right),
\]
with $\mathcal{N}=B\cup R$.

At each simulator tick $t$, agent $i$ receives a local observation
$o_{i,t}\sim O_i(\cdot\mid s_t)$ and may request a temporally extended
\emph{macro-action} together with a target. The macro vocabulary is
$\{\texttt{GO\_TO},\texttt{GRAB\_MINE},\texttt{GET\_FLAG},\texttt{PLACE\_MINE},\texttt{GO\_HOME}\}$.
For reproducibility, the joint action is represented as
$\mathrm{MultiDiscrete}([5,50,5,50])$ (per agent: a $5$-way macro and a
$50$-way target). The joint action induces $s_{t+1}\sim P(\cdot\mid s_t,a_t)$.

Each team's flag is anchored at its home base; a carrier that returns the enemy
flag alive scores a capture. Mines are stationary hazards that create
area-denial and ambush dynamics.

\subsection{Asynchronous macro-action decision points}

An accepted macro-action is held for a macro-specific number of ticks before the
agent may choose a new one
($\texttt{GO\_TO}=4$, $\texttt{GRAB\_MINE}=3$, $\texttt{GET\_FLAG}=4$,
$\texttt{PLACE\_MINE}=2$, $\texttt{GO\_HOME}=4$). Agent $i$ carries a countdown
$c_{i,t}$; a newly requested macro is accepted only when
\[
d_t^i=\mathbb{I}\!\left[c_{i,t}\le 0\right].
\]
When $d_t^i=0$, the legal-action mask collapses to the already-committed macro,
so the network output is not an executable decision. Because hold lengths differ
by macro and agents commit independently, decision boundaries are asynchronous.
For policy evaluation and supervision, the task is therefore better viewed as a
semi-Markov decision process over commitment decisions than as a flat per-tick
Markov decision process.

\subsection{Opponent regimes}

Red is drawn from one of two scripted defensive philosophies that we call
\emph{opponent regimes}:
\begin{itemize}
  \item \textbf{Opponent Regime~A:} broad / passive zone defense.
  \item \textbf{Opponent Regime~B:} compact / reactive defense.
\end{itemize}
The two regimes are designed to elicit contrasting strategic responses from
Blue rather than represent two difficulty levels of the same opponent behavior.
Implementation identifiers used for reproducibility (configuration overlay
\textbf{OP6} for Regime~A; canonical \textbf{OP7} for Regime~B), together with
frozen behavior-tree parameters such as
$\texttt{defender\_zone\_frac}$ and $\texttt{threat\_radius}$, are recorded in
the experiment configuration and are not repeated as conceptual labels in the
main narrative.

We do not assume a priori that the two regimes require different Blue
strategies. An \emph{opponent-conditioned strategy-preference test} measures
whether a GUARD-biased response is preferred over a BREACH-biased response under
Regime~A, and the reverse contrast under Regime~B, with pass condition
$\mathrm{LCB}_{95}>0$ on each contrast. This empirical test establishes whether
the regimes induce complementary strategic demand. If one Blue policy were
simply better against both regimes, a single shared behavior would be
appropriate rather than a specialization problem.

\subsection{Strategy-conditioned policies}

Let $\pi_\theta(a\mid o,z)$ denote a shared multi-robot policy conditioned on a
discrete \emph{strategy code} $z\in\{z_0,z_1\}$. The strategy code selects which
opponent-conditioned behavior the shared policy should express. The intended
deployment mapping is $z_0\mapsto\text{Regime~A}$ and $z_1\mapsto\text{Regime~B}$;
once fixed, this assignment is never flipped. Let $V(z,p)$ denote the win rate
when deploying with strategy code $z$ against opponent regime
$p\in\{A,B\}$, estimated over a held-out seed block.

Independently trained \emph{strategy-specific expert policies} $\pi_A$ and
$\pi_B$ are the source behaviors: each is trained against one regime with no
shared parameters and no strategy-code mechanism. Their bit-exact $z$-dispatch
wrapper is Share-0 (Section~\ref{sec:sharing}). A single \emph{generalist}
policy $\pi_G$ trained on the mixed opponent distribution without a strategy
code provides a single-policy baseline.

\subsection{Strategy-specialization criterion}
\label{sec:specialization}

We require more than behavioral diversity. For two strategy codes to represent
meaningful specialization, each code must provide higher task success than the
alternative code in its intended opponent regime. A policy is considered
strategy-specialized only if each strategy code significantly outperforms the
alternative code in the opponent regime for which it is intended.

Formally, define
\[
\Delta_A = V(z_0,A)-V(z_1,A),
\qquad
\Delta_B = V(z_1,B)-V(z_0,B).
\]
The \emph{strategy-specialization criterion} (equivalently, the two-sided
\emph{specialization crossover} pattern) requires
\[
\mathrm{LCB}_{95}(\Delta_A)>0
\quad\wedge\quad
\mathrm{LCB}_{95}(\Delta_B)>0,
\]
evaluated by a paired percentile bootstrap over the held-out seeds with
$n_{\mathrm{boot}}=20{,}000$, $\alpha=0.05$, and the seed as the resampling
unit. The binary win indicator is the primary outcome because it is the
deployment quantity used by every specialization evaluation.

Independent-expert specialization is reported on a held-out $64$-seed block.
All parameter-sharing comparisons use the matched $128$-seed paired protocol of
Section~\ref{sec:matched}.

\subsection{Scope of claims}

Primary claims apply to the nominal $2$-vs-$2$ environment, the two opponent
regimes above, the specific training recipe, and the fixed evaluation
implementation. Scaling to $4$-vs-$4$ and $6$-vs-$6$, and robustness
perturbations such as localization noise, motion uncertainty, and control delay,
are external-validity experiments. Until completed, they are not used as evidence
for the primary sharing conclusions and are not used to select the primary
architecture.

%======================================================================
